% Exercises for Chapter 4 (The integral of non-negative functions). Statements only:
% solutions are kept privately. Throughout, (X, A, mu) is a measure space, m is the
% Borel-Lebesgue measure on B(R), and f_n are the dyadic approximants of Chapter 3.

\begin{exercise}[Rules for the integral of simple functions]\label{ex:in-simple-rules}
\easy This exercise revisits the properties of the integral of simple functions that the notes
leave to the reader (\noteref{prop:in-simple-properties}). You may use
\noteref{def:in-simple-integral}, \noteref{prop:in-simple-well-defined}, properties (1)--(3) of
\noteref{prop:in-simple-properties} and \noteref{prop:mf-common-partition}, with the convention
$0\cdot\infty=0$.
\begin{enumerate}[label=(\alph*)]
  \item Let $A_1,\dots,A_k\in\mathcal A$ be arbitrary measurable sets, not necessarily disjoint
    and not necessarily covering $X$, and let $a_1,\dots,a_k\in[0,+\infty]$. Show that
    $\varphi=\sum_{i=1}^ka_i\1_{A_i}$ belongs to $\mathcal S^+(X)$ and that
    $\int_X\varphi\,\dd\mu=\sum_{i=1}^ka_i\,\mu(A_i)$.
  \item Let $\varphi,\psi\in\mathcal S^+(X)$ with $\varphi\le\psi$. Construct $\eta\in\mathcal
    S^+(X)$ with $\psi=\varphi+\eta$ (where $\varphi=+\infty$ the difference $\psi-\varphi$ is not
    defined, so choose $\eta$ there). Deduce property (4), monotonicity, from property (3).
  \item Let $E,F\in\mathcal A$ be disjoint. Show that
    $\int_{E\cup F}\varphi\,\dd\mu=\int_E\varphi\,\dd\mu+\int_F\varphi\,\dd\mu$, and deduce property
    (5).
  \item Show that if $\int_X\varphi\,\dd\mu<+\infty$, then $\mu(\{\varphi=+\infty\})=0$ and
    $\mu(\{\varphi>0\})<+\infty$, and that the converse holds. Show that the converse fails for
    measurable functions that are not simple: on $((0,1),\mathcal B((0,1)),m)$, find $f$ with
    finite values and $\int f\,\dd m=+\infty$.
\end{enumerate}
\end{exercise}

\begin{exercise}[Approximation from below, not from above]\label{ex:in-from-below}
\medium The integral of \noteref{def:in-integral} approximates $f$ from below by simple
functions (and, as the proof of \noteref{thm:in-monotone-convergence} shows, finite-valued ones
suffice). This exercise shows that approximating from above would not work. Let
$f:X\to[0,+\infty]$ be measurable. You may use \noteref{sec:in-simple},
\noteref{prop:in-properties}, \noteref{cor:in-series} and \noteref{rem:mf-dyadic}.
\begin{enumerate}[label=(\alph*)]
  \item Let $\int^*f\,\dd\mu=\inf\{\int_X\psi\,\dd\mu:\psi\in\mathcal S^+(X),\ \psi\ge f\}$, the
    \emph{upper integral}. Show that $\int_Xf\,\dd\mu\le\int^*f\,\dd\mu$.
  \item Suppose that $\mu(X)<+\infty$ and that $f$ is bounded. Show that
    $\int^*f\,\dd\mu=\int_Xf\,\dd\mu$. \emph{Hint:} compare $f$ with $f_n$ and $f_n+2^{-n}$.
  \item On $((0,1),\mathcal B((0,1)),m)$ let $f(x)=x^{-1/2}$. Show that $f$ is integrable, by
    comparing it with $\sum_{j\ge0}2^{(j+1)/2}\1_{[2^{-j-1},2^{-j})}$, and that
    $\int^*f\,\dd m=+\infty$.
  \item On $(\mathbb R,\mathcal B(\mathbb R),m)$ let $f(x)=e^{-|x|}$. Show that $f$ is
    integrable and bounded, but $\int^*f\,\dd m=+\infty$. Which hypothesis of (b) fails in (c),
    and which in (d)?
\end{enumerate}
\end{exercise}

\begin{exercise}[The layer-cake formula]\label{ex:in-layer-cake}
\hard Let $f:X\to[0,+\infty]$ be measurable, let $f_n$ be its dyadic approximants
(\noteref{eq:mf-dyadic}), and let $G(t)=\mu(\{f\ge t\})$ for $t>0$. You may use the results of
the chapter.
\begin{enumerate}[label=(\alph*)]
  \item Show that $f_n=2^{-n}\sum_{k=1}^{n2^n}\1_{\{f\ge k2^{-n}\}}$, and deduce that
    $\int_Xf_n\,\dd\mu=2^{-n}\sum_{k=1}^{n2^n}G(k2^{-n})$.
  \item Deduce that $\int_Xf\,\dd\mu=\lim_n2^{-n}\sum_{k=1}^{n2^n}G(k2^{-n})$.
  \item Suppose that $G(t)<+\infty$ for every $t>0$. Since $G$ is non-increasing, it is Riemann
    integrable on every interval $[a,b]\subseteq(0,+\infty)$. Show that
    \[
      \int_Xf\,\dd\mu=\int_0^{+\infty}G(t)\,\dd t=\lim_{a\downarrow0,\ b\uparrow+\infty}\int_a^bG(t)\,\dd t,
    \]
    an improper Riemann integral, possibly $+\infty$. \emph{Hint:} bound $2^{-n}G(k2^{-n})$ above
    and below by integrals of $G$ over neighbouring intervals. What happens if
    $G(t_0)=+\infty$ for some $t_0>0$?
  \item Let $(\Omega,\mathcal F,\mathbb P)$ be a probability space and $Y:\Omega\to[0,+\infty)$ a
    random variable. Show that
    $\int_\Omega Y\,\dd\mathbb P=\int_0^{+\infty}\mathbb P(Y\ge t)\,\dd t=\int_0^{+\infty}\mathbb P(Y>t)\,\dd t$,
    and that $\int_\Omega Y\,\dd\mathbb P=\sum_{k\ge1}\mathbb P(Y\ge k)$ if $Y$ takes values in
    $\mathbb N$.
  \item Compute $\int_{(0,1)}x^{-p}\,\dd m(x)$ for every $p>0$.
\end{enumerate}
\end{exercise}

\begin{exercise}[How fast does $\int f_n$ approach $\int f$?]\label{ex:in-dyadic-rates}
\medium This is a computational exercise on the dyadic approximants $f_n$
(\noteref{eq:mf-dyadic}) and on the rate in \noteref{rem:in-dyadic}. On
$(\mathbb R,\mathcal B(\mathbb R),m)$ let $\delta_n=\int f\,\dd m-\int f_n\,\dd m$. You may use
\cref{ex:in-layer-cake}, in particular
$\int f_n\,\dd\mu=2^{-n}\sum_{k=1}^{n2^n}\mu(\{f\ge k2^{-n}\})$.
\begin{enumerate}[label=(\alph*)]
  \item Write a vectorised function that computes $\int f_n\,\dd\mu$ from the function
    $t\mapsto\mu(\{f\ge t\})$, summing in chunks so that $n=20$ fits in memory. Check it on a
    finite space with a thousand points against $\sum_xf_n(x)\,\mu(\{x\})$. Prove that for
    $f(x)=x$ on $[0,1]$ (and $0$ elsewhere) $\int f_n\,\dd m=\frac12-2^{-n-1}$.
  \item For $f(x)=x^2$ on $[0,1]$, $f(x)=x^{-1/2}$ on $(0,1)$ (each $0$ elsewhere) and
    $f(x)=1/(1+x^2)$ on $\mathbb R$, find $m(\{f\ge t\})$ and $\int f\,\dd m$, and tabulate
    $2^{n+1}\delta_n$, $n\,\delta_n$ and $2^{n/2}\delta_n$ for $n=1,\dots,18$. Which of the three
    settles down for each function?
  \item Write $f-f_n=r_n+c_n+s_n$, where $r_n=(f-f_n)\1_{\{2^{-n}\le f<n\}}$ is the rounding
    error, $c_n=(f-n)\1_{\{f\ge n\}}$ is the cut, and $s_n=f\1_{\{f<2^{-n}\}}$ collects the small
    values that $f_n$ discards. Show that $\int r_n\,\dd m\le2^{-n}m(\{2^{-n}\le f<n\})$, compute
    $\int c_n\,\dd m$ and $\int s_n\,\dd m$ exactly for the three functions, and explain the rates
    found in (b).
  \item For $f(x)=1/(1+x^2)$, the rule of thumb ``the rounding error is half a step on average''
    predicts $\delta_n\approx3\cdot2^{-n/2}$. Show that in fact
    $2^{n/2}\delta_n\to-2\zeta(\frac12)\approx2.9207$, where $\zeta(\frac12)\approx-1.46035$ is the
    constant in $\sum_{k=1}^Nk^{-1/2}=2\sqrt N+\zeta(\frac12)+O(N^{-1/2})$. \emph{Hint:}
    $\sqrt{1-u}=1-\sqrt u\,F(u)$ with $F(u)=\sqrt u/(1+\sqrt{1-u})$, which is increasing on
    $[0,1]$; compare a Riemann sum of $F$ with $\int_0^1F$. Why does the rule of thumb fail here,
    although it works for $x^2$?
\end{enumerate}
\end{exercise}

\begin{exercise}[Markov's inequality]\label{ex:in-markov}
\easy Let $f:X\to[0,+\infty]$ be measurable. You may use \noteref{prop:in-properties},
\noteref{prop:in-additivity} and \noteref{prop:in-ae}.
\begin{enumerate}[label=(\alph*)]
  \item Show that $\mu(\{f\ge t\})\le\frac1t\int_Xf\,\dd\mu$ for every $t\in(0,+\infty)$.
  \item Deduce, in one line each, item (1) of \noteref{prop:in-ae} and the implication
    ``$\int_Xf\,\dd\mu=0\Rightarrow f=0$ a.e.'' of its item (2).
  \item Show that if $f$ is integrable, then $\{f>0\}$ is a countable union of measurable sets
    of finite measure.
  \item Suppose that $f$ is integrable and $t>0$. Show that equality holds in (a) if and only if
    $f(x)\in\{0,t\}$ for almost every $x$.
  \item Let $(\Omega,\mathcal F,\mathbb P)$ be a probability space and $Y:\Omega\to\mathbb R$ a
    random variable. Show \emph{Chebyshev's inequality}
    $\mathbb P(|Y|\ge t)\le t^{-2}\int_\Omega Y^2\,\dd\mathbb P$ for $t>0$.
\end{enumerate}
\end{exercise}

\begin{exercise}[Series of non-negative functions at work]\label{ex:in-series-uses}
\medium You may use \noteref{cor:in-series}, \noteref{prop:in-ae} and the results of
\cref{ch:02-measures}.
\begin{enumerate}[label=(\alph*)]
  \item Let $(A_n)_{n\ge1}$ be measurable sets with $\sum_n\mu(A_n)<+\infty$, and let
    $N=\sum_{n\ge1}\1_{A_n}$, so that $N(x)$ counts the sets $A_n$ that contain $x$. Show that
    $\int_XN\,\dd\mu=\sum_n\mu(A_n)$ and $\limsup_nA_n=\{N=+\infty\}$, and deduce the first
    Borel--Cantelli lemma (\noteref{lem:ms-borel-cantelli-1}).
  \item Let $f_n\ge0$ be measurable with $\sum_n\int_Xf_n\,\dd\mu<+\infty$. Show that $f_n\to0$
    a.e.
  \item For $q\in[0,1]$ let $h_q(x)=|x-q|^{-1/2}$ for $x\ne q$ and $h_q(q)=+\infty$. Show that
    $\int_{[0,1]}h_q\,\dd m\le2(1+\sqrt2)$, by comparing $h_q$ on $[0,1]$ with
    $\sum_{j\ge0}2^{(j+1)/2}\1_{\{2^{-j-1}\le|x-q|<2^{-j}\}}$.
  \item Let $(q_k)_{k\ge1}$ be an enumeration of $\mathbb Q\cap[0,1]$ and
    $g=\sum_{k\ge1}2^{-k}h_{q_k}$ on $[0,1]$. Show that $g$ is integrable on $[0,1]$, hence finite
    a.e., although $g=+\infty$ at every rational point and $g$ is unbounded on every
    non-degenerate interval.
\end{enumerate}
\end{exercise}

\begin{exercise}[Small sets carry a small integral]\label{ex:in-abs-continuity}
\medium Let $f:X\to[0,+\infty]$ be measurable and integrable. You may use the results of the
chapter.
\begin{enumerate}[label=(\alph*)]
  \item Show that $\int_{\{f>M\}}f\,\dd\mu\to0$ as $M\to+\infty$. \emph{Hint:} apply monotone
    convergence to $f\1_{\{f\le M\}}$, and use \noteref{prop:in-ae}.
  \item Show that for every $\varepsilon>0$ there is $\delta>0$ such that
    $\int_Ef\,\dd\mu<\varepsilon$ for every $E\in\mathcal A$ with $\mu(E)<\delta$.
  \item On $((0,1),\mathcal B((0,1)),m)$, show that $f(x)=1/x$ has $\int_{(0,\delta)}f\,\dd m=+\infty$
    for every $\delta\in(0,1)$: (b) fails without integrability, even for a function with finite
    values.
  \item Show that $M\mu(\{f>M\})\to0$ as $M\to+\infty$, which improves Markov's inequality
    (\cref{ex:in-markov}) for large $M$.
\end{enumerate}
\end{exercise}

\begin{exercise}[Decreasing sequences]\label{ex:in-decreasing}
\medium This exercise is partly computational. Let $(f_n)_{n\ge1}$ be non-negative measurable
functions with $f_n\ge f_{n+1}$, and $f=\lim_nf_n$. You may use the results of the chapter and
\cref{ex:in-layer-cake}.
\begin{enumerate}[label=(\alph*)]
  \item Show that if $\int_Xf_1\,\dd\mu<+\infty$, then $\int_Xf_n\,\dd\mu\to\int_Xf\,\dd\mu$.
    \emph{Hint:} $f_1<+\infty$ a.e.; apply monotone convergence to $f_1-f_n$ on
    $\{f_1<+\infty\}$, and use additivity.
  \item Show that the hypothesis $\int f_1\,\dd\mu<+\infty$ cannot be dropped: on
    $(\mathbb R,\mathcal B(\mathbb R),m)$ take $f_n=\1_{[n,+\infty)}$, and on
    $((0,1),\mathcal B((0,1)),m)$ take $f_n(x)=x^{-1}\1_{(0,1/n)}(x)$. In both cases
    $f_n\searrow0$ but $\int f_n\,\dd m=+\infty$ for every $n$.
  \item For $f_n(x)=x^{-1}\1_{(0,1/n)}(x)$, compute the integrals of the dyadic approximants
    $\varphi_k$ of $f_n$, for $k=1,\dots,20$ and $n=1,2,4$, from $m(\{f_n\ge t\})$. Show that they
    equal $k/n$ for $k\le n$ and $1+\sum_{j=n2^k+1}^{k2^k}\frac1j$ for $k\ge n$, so that they grow
    like $1+\log(k/n)$. How does this show that $\int f_n\,\dd m=+\infty$?
  \item Repeat (c) for $g_n(x)=x^{-1/2}\1_{(0,1/n)}(x)$. Show that $\int g_n\,\dd m=2/\sqrt n$,
    which tends to $0$ as (a) predicts, and that, once $k\ge\sqrt n$, the dyadic integrals fall
    short of it by $\frac1k$, up to a rounding error below $2^{-k}/n$.
\end{enumerate}
\end{exercise}

\begin{exercise}[When is Fatou's lemma strict?]\label{ex:in-fatou-strict}
\medium This exercise is partly computational. You may use the results of the chapter.
\begin{enumerate}[label=(\alph*)]
  \item Show that Fatou's lemma and monotonicity together imply the monotone convergence
    theorem.
  \item Let $X=\{0,\dots,N-1\}$ with $\mu(\{x\})=w_x>0$, and let $f_n=h_{n\bmod p}$ be a periodic
    sequence of functions $X\to[0,+\infty)$, with period $p$. Show that
    $\liminf_nf_n=\min_jh_j$ pointwise and $\liminf_n\int f_n\,\dd\mu=\min_j\int h_j\,\dd\mu$, and
    that equality holds in Fatou's lemma if and only if some $h_j$ equals $\min_ih_i$ at every
    point.
  \item Draw $10^4$ random examples with $p=3$ and $N=6$, with independent values $h_j(x)$
    uniform in $\{0,1,2,3\}$ and weights $w_x$ uniform in $(0,1)$. Check Fatou's lemma and the
    criterion of (b) in each case, and compare the fraction of strict inequalities with its
    exact probability, which you should compute from (b). Find by hand an example with $p=N=3$
    and $w_x=1$ in which the two sides are $0$ and $2$.
  \item On $(\mathbb R,\mathcal B(\mathbb R),m)$ let $f_n=\1_{[0,1)}+\1_{[n,n+1)}$. Compute, by hand
    and with step functions, $\int f_n\,\dd m$ and $\int\inf_{n\le k\le K}f_k\,\dd m$ for
    $K=50$, and show that Fatou's lemma reads $1<2$.
  \item Show that Fatou's lemma becomes false if $\liminf$ is replaced by $\limsup$ on both
    sides, and that the reverse inequality $\limsup_n\int f_n\,\dd\mu\le\int\limsup_nf_n\,\dd\mu$
    can fail too. \emph{Hint:} $f_n=\1_{[0,1)}$ for even $n$ and $\1_{[1,2)}$ for odd $n$; the
    sliding bump.
\end{enumerate}
\end{exercise}
